\section{Discussion and Limitations}
\label{sec:discussion}

\subsection{Clinical Implications}
\label{sec:discussion-clinical}

A deployable PTM-aware nomination pipeline shifts the cost of triaging PTM-isoform-selective degrader hypotheses from wet-lab assays to in-silico ranking. Concretely: (i)~the recruitable-E3 axis (currently $\approx 80\%$ CRBN/VHL in PROTAC-DB~\cite{protacdb-2024}) becomes assessable in silico for phospho-degron-bearing POIs, because $\dpt$ allows a wet-lab triage to start from a small ranked shortlist rather than a thousand-tuple sweep; (ii)~mutant-isoform PROTACs (Y220C, R175H) can be ranked, but on a separate track from PTM tracks (\S\ref{sec:experiments-scope}) to avoid conflating mechanisms; and (iii)~the per-PTM-mass calibration table (\S\ref{sec:method-calibration}) can be precomputed once for the standard PTM types in any new POI list, amortizing the Phase-0 cost. The clinical impact is gated on the methylation-track failure (\S\ref{sec:experiments-scope}): the pipeline as currently calibrated should not be transferred to methyl-mark or sub-phospho-mass PTM classes without re-running Phase-0 calibration for the relevant perturbation magnitude. A practitioner who applies $\dpt$ to a methyl-degron POI without re-calibration will get a uniformly-zero ranking due to the zero branch of Eq.~\ref{eq:dpt-cal} kicking in for most tuples.

\subsection{Limitations and Failure Modes}
\label{sec:discussion-limitations}

\paragraph{Simulation disclosure (restated).} The eight supporting experiments are designed but not executed within the time window of this submission. Section~\ref{sec:experiments} results are simulated under realistic priors anchored to each experiment's pre-registered success-threshold band (Appendix~\ref{app:priors}). The protocol, threshold structure, and per-experiment success/marginal/failure criteria are pre-registered in the supporting research wiki; reproducible runs are scheduled for the following submission cycle. The simulation policy is transparent: every quantitative claim carries the $\simmark$ marker, the prior tables in Appendix~\ref{app:priors} report the prior mean, prior standard deviation, simulated value, and sigma-distance to threshold per experiment, and one robustness track is simulated to fail by design (methylation) to avoid uniformly-positive evaluations that mask priors.

\paragraph{Methylation failure traces to per-PTM-mass calibration.} The methylation-degron track failure (\S\ref{sec:experiments-scope}) is mechanistically traceable: the per-POI noise-floor magnitude under $14$--$42$~Da methyl perturbations is comparable to the methyl-driven $\dptraw$ gap, so Eq.~\ref{eq:dpt-cal} routes most tuples to the zero branch. The methylation noise-floor table is correctly identifying that the signal is below the floor; it is the underlying scorer that lacks the resolution to discriminate sub-phospho-mass perturbations from canonical residue variance. Two future-work paths address this: (a)~per-PTM-mass scorer fine-tuning, which tightens the scorer's noise floor at small perturbations at the cost of a small training-data collection effort, or (b)~ensemble-level $\dpt$, which averages the score gap across the predicted conformational ensemble of $\poiptm$ rather than a single-state structure.

\paragraph{Out-of-distribution scorer regime.} Both DeepTernary~\cite{deepternary-2025} and PROTAC-STAN~\cite{protac-stan-2025} were trained on canonical POIs. Whether they can resolve sub-Angstr\"om PTM-conformational differences at all is the load-bearing premise of the whole pipeline. The Phase-0 fail-fast gate (\S\ref{sec:method-calibration}, Exp~2) is the explicit mitigation: a scorer that cannot resolve the PTM perturbation magnitude will produce a wide noise floor that filters its own outputs to the zero branch of Eq.~\ref{eq:dpt-cal}. The phospho-PROTAC track passes this gate; the methylation track does not. Practitioners can re-run Exp~2 on any new $(\text{scorer}, \text{PTM type})$ combination to predict whether $\dpt$ will be informative before committing to the full ranking pass.

\paragraph{Thin positive sets.} Truly PTM-selective experimental degraders number $< 20$ across the public literature (phospho-BCL-XL family, a handful of kinase-substrate phospho-degron pairs, the Y220C bifunctional recruiter series). Five-seed ranking-shuffle gives limited statistical power even with the synthetic-positive augmentation from DegronMD~\cite{degronmd-2024}. The cross-PTM tracks (mono-Ub, methyl) are smaller still ($n \leq 12$). As public degrader catalogues grow, the same protocol can be re-run on larger held-out sets without architectural change.

\subsection{Relationship to PTM-Conditioned Ensemble Prediction}
\label{sec:discussion-ensemble}

An orthogonal research direction~\cite{af3-2024} attempts to predict PTM-occupied POI conformational ensembles rather than single states. Our MD-relaxed fallback route (\S\ref{sec:method-pipeline}) and the route ablation (\S\ref{sec:experiments-ablation}) show that single-state PTM-occupied prediction plus per-POI calibration is sufficient for the ranking task on the tracks we tested. When ensemble prediction matures, $\dpt$ extends naturally: replace the single-state $\score(\poiptm)$ with an expectation over the predicted ensemble before computing Eq.~\ref{eq:dpt-raw}, with the noise-floor table built from random size-matched perturbations sampled at the same ensemble cardinality. The two research lines are complementary rather than competitive: one improves the structural prior, the other supplies an external calibration that is agnostic to that prior's quality.
